\frontmatter%前言
\chapter*{前言}
\addcontentsline{toc}{chapter}{前言}
信号与系统是电子信息等相关专业本科生的重要基础课，是后续课程“数字信号处理”“通信原理”“自动控制”“随机信号处理”“数字图像处理”“现代信号处理”等课程的基础，以确定信号通过线性时不变系统的过程为核心，综合运用多门数学课程的知识，尤其是傅里叶分析。

然而，国内外的信号与系统课程对于傅里叶变换的定义有所不同，国内的课程中所谓“频域分析”实际上指的是角频率域，而这种定义上的不同会导致诸多公式的内核相同但形式不同，使学习者感到困惑。我自学过斯坦福大学公开课EE261，同样面临着这个问题，所以希望尝试这种新的笔记形式，梳理信号与系统课程的理论部分，同时也使本书能够作为同学们学习信号与系统时的参考书。我们将看到，同时给出两种傅里叶变换，不仅方便我们对比、查找公式，理解傅里叶变换的本质，还能在必要时转向其中一种傅里叶变换，使得讨论更加简洁明快。

同时，我注意到现在的信号与系统教材中对很多数学上的内容的处理过于简单化，以至于同学们对一些公式感到不解，我希望结合《信号与系统》教材和实变函数、复变函数和傅里叶分析等教材，对这样的内容做一些补充和深化。对于其中比较复杂的内容，我会用“*”进行标注，以提示读者可以选择性地阅读。

本书将不会涉及信号与系统实践部分，对一些语文性质的概念从简处理，希望了解这些内容的读者可以自行查看信号与系统教材。对于两种不同的傅里叶变换，将并排给出两种公式，左侧为频率版本，右侧为角频率版本。为了逻辑的完整性和内容的精确性，本书会引入一些远超工科教学范围的内容，力求朴素平和却又不失深度，其中所涉及的技术细节不在我们的关注范围，读者也可以自行选择是否阅读这些内容。本书默认读者已经学习过一到两学期的数学分析（或高等数学/微积分）以及高等代数（或线性代数）课程，不过只要理解一元微积分，阅读前四章及就不会有问题；只要掌握复数运算、理解曲线积分，阅读后两章就不会有问题。

限于作者水平，书中难免存在不足之处，希望读者可以提出宝贵的意见或建议，可以反馈至我的电子邮箱：2024212622@bupt.cn。

{
\raggedleft{}
北京市十一学校2024届毕业生\\
严子竣\\
2025年秋\\
}

\chapter*{致谢}
\addcontentsline{toc}{chapter}{致谢}
